% =============================================================================
% Appendix B: Use Cases
% NexaLearn Graduation Project — Cairo University, Faculty of Engineering
% =============================================================================

\chapter{Use Cases}
\label{app:usecases}

This appendix documents the ten primary use cases of the NexaLearn platform. Each use case is presented in a structured format with an identifier, name, actors, preconditions, main flow, postconditions, and exception conditions.

% ---------------------------------------------------------------------------
\section{Use Case Descriptions}

% UC-01
\subsection{UC-01: Student Signup and Email Verification}

\begin{longtable}{p{3cm}p{10cm}}
  \toprule
  \textbf{Element} & \textbf{Description} \\
  \midrule
  \endhead
  \midrule
  \endfoot
  \textbf{ID} & UC-01 \\
  \textbf{Name} & Student Signup and Email Verification \\
  \textbf{Actors} & Unauthenticated student \\
  \textbf{Preconditions} & The student has an email address and a password that meets the platform's password policy (minimum 8 characters, at least one uppercase letter and one digit). \\
  \textbf{Main Flow} &
  \begin{enumerate}
    \item The student submits a registration request with full name, email address, and password via \texttt{POST /auth/register}.
    \item The system validates the input format (email syntax, password strength) and checks that the email address is not already registered.
    \item The system creates a \texttt{User} aggregate in \texttt{UNVERIFIED} status, hashes the password using bcrypt (cost factor 12), and stores the record in PostgreSQL.
    \item The system generates a cryptographically random email verification token (64 hex characters), stores its SHA-256 hash in Redis with a 24-hour TTL, and sends an email via Resend containing a verification link: \texttt{https://nexalearn.app/auth/verify?token=<raw\_token>}.
    \item The student clicks the verification link. The system hashes the raw token and retrieves the corresponding user ID from Redis.
    \item The system updates the \texttt{User} aggregate to \texttt{ACTIVE} status and returns a success response.
  \end{enumerate} \\
  \textbf{Postconditions} & The user is registered and verified. The user can log in using email and password. \\
  \textbf{Exceptions} &
  \begin{itemize}
    \item \textbf{E01: Duplicate email.} If the email is already registered, the system returns HTTP 409 Conflict with the message "An account with this email already exists."
    \item \textbf{E02: Invalid token.} If the verification token is expired or malformed, the system returns HTTP 400 Bad Request.
    \item \textbf{E03: Email delivery failure.} If the Resend API returns a non-2xx response, the user is created but remains \texttt{UNVERIFIED}. A resend endpoint (\texttt{POST /auth/resend-verification}) allows the user to request a new token.
  \end{itemize} \\
  \bottomrule
\end{longtable}

% UC-02
\subsection{UC-02: Instructor Course Creation and Publishing}

\begin{longtable}{p{3cm}p{10cm}}
  \toprule
  \textbf{Element} & \textbf{Description} \\
  \midrule
  \endhead
  \midrule
  \endfoot
  \textbf{ID} & UC-02 \\
  \textbf{Name} & Instructor Course Creation and Publishing \\
  \textbf{Actors} & Authenticated instructor \\
  \textbf{Preconditions} & The instructor is logged in with a valid access token. The instructor's role is \texttt{INSTRUCTOR}. \\
  \textbf{Main Flow} &
  \begin{enumerate}
    \item The instructor submits course metadata (title, description, category, difficulty level, price, cover image) via \texttt{POST /courses}. The course is created in \texttt{DRAFT} status.
    \item The instructor adds modules to the course via \texttt{POST /courses/:id/modules}. Each module has a title, description, and ordinal position.
    \item The instructor adds lessons to each module via \texttt{POST /modules/:id/lessons}. Each lesson has a title, content (HTML/Markdown), optional video asset, and optional quiz.
    \item The instructor uploads a video for a lesson via the Mux direct upload flow (see UC-04).
    \item The instructor submits a publish request via \texttt{POST /courses/:id/publish}.
    \item The system validates the course completeness: at least one module required, each module must have at least one lesson, all lessons must have either text content or a video asset.
    \item The system transitions the \texttt{Course} aggregate from \texttt{DRAFT} to \texttt{PUBLISHED} and publishes a \texttt{CoursePublished} domain event through the outbox.
  \end{enumerate} \\
  \textbf{Postconditions} & The course is visible in the course catalogue. Students can browse the course and enroll. \\
  \textbf{Exceptions} &
  \begin{itemize}
    \item \textbf{E01: Incomplete course.} If any completeness validation fails, the system returns HTTP 422 Unprocessable Entity with a list of missing requirements.
    \item \textbf{E02: Unauthorised.} If the authenticated user's role is not \texttt{INSTRUCTOR} or \texttt{ADMIN}, the system returns HTTP 403 Forbidden.
  \end{itemize} \\
  \bottomrule
\end{longtable}

% UC-03
\subsection{UC-03: Student Enrollment in Paid Course}

\begin{longtable}{p{3cm}p{10cm}}
  \toprule
  \textbf{Element} & \textbf{Description} \\
  \midrule
  \endhead
  \midrule
  \endfoot
  \textbf{ID} & UC-03 \\
  \textbf{Name} & Student Enrollment in Paid Course (Stripe Payment) \\
  \textbf{Actors} & Authenticated student, Stripe payment gateway \\
  \textbf{Preconditions} & The student is logged in. The course is in \texttt{PUBLISHED} status and has a non-zero price. \\
  \textbf{Main Flow} &
  \begin{enumerate}
    \item The student requests enrollment via \texttt{POST /enrollments/:courseId}.
    \item The system checks that the student is not already enrolled (unique constraint on \texttt{(user\_id, course\_id, status)} where \texttt{status != 'CANCELLED'}).
    \item The system creates a \texttt{PaymentIntent} record in \texttt{PENDING} status and requests a payment intent from Stripe via \texttt{stripe.paymentIntents.create()} with the course price and currency.
    \item The system returns the \texttt{client\_secret} to the student's frontend.
    \item The student completes the payment on the frontend using Stripe Elements or Stripe Checkout.
    \item Stripe delivers a \texttt{payment\_intent.succeeded} webhook event to \texttt{POST /webhooks/stripe}.
    \item The system verifies the webhook signature, checks idempotency via the \texttt{ProjectionEventLedger}, updates the \texttt{PaymentIntent} to \texttt{SUCCEEDED}, creates an \texttt{Enrollment} aggregate in \texttt{ACTIVE} status, and publishes an \texttt{EnrollmentActivated} event.
  \end{enumerate} \\
  \textbf{Postconditions} & The student is enrolled in the course with \texttt{ACTIVE} status. The student can access all lessons. \\
  \textbf{Exceptions} &
  \begin{itemize}
    \item \textbf{E01: Already enrolled.} If the student is already enrolled, the system returns HTTP 409 Conflict.
    \item \textbf{E02: Payment failed.} If Stripe reports a payment failure, the \texttt{PaymentIntent} transitions to \texttt{FAILED} and the enrollment is created in \texttt{PENDING\_PAYMENT} status with a link to retry payment.
    \item \textbf{E03: Duplicate webhook.} If the webhook is delivered multiple times, the \texttt{ProjectionEventLedger} unique constraint prevents duplicate processing.
  \end{itemize} \\
  \bottomrule
\end{longtable}

% UC-04
\subsection{UC-04: Video Upload and AI Quiz Generation}

\begin{longtable}{p{3cm}p{10cm}}
  \toprule
  \textbf{Element} & \textbf{Description} \\
  \midrule
  \endhead
  \midrule
  \endfoot
  \textbf{ID} & UC-04 \\
  \textbf{Name} & Video Upload and AI Quiz Generation \\
  \textbf{Actors} & Authenticated instructor, Mux Video API, OpenAI Whisper / LLM \\
  \textbf{Preconditions} & The instructor is logged in. A lesson exists in the course. \\
  \textbf{Main Flow} &
  \begin{enumerate}
    \item The instructor requests a Mux direct upload URL via \texttt{POST /media/upload}. The system calls Mux's \texttt{/video/uploads} API and returns the upload URL to the frontend.
    \item The frontend uploads the video file directly to Mux. Mux fires \texttt{video.asset.created} webhook upon upload completion.
    \item The system creates a \texttt{VideoAsset} aggregate in \texttt{UPLOADED} status.
    \item Mux transcodes the video and fires \texttt{video.asset.ready} (HLS available) and later \texttt{video.asset.mp4\_rendition.ready}.
    \item When the asset reaches \texttt{MP4\_AVAILABLE} status, the system creates a \texttt{TranscriptionJob} in \texttt{PENDING} status and sends the MP4 download URL to the AI pipeline with an HMAC-signed callback URL.
    \item The AI pipeline (Whisper) transcribes the audio and POSTs the transcript back to \texttt{POST /webhooks/ai-pipeline/transcription}. The system verifies the HMAC signature, stores the transcript, and creates a \texttt{QuizGenerationJob} in \texttt{PENDING} status.
    \item The AI pipeline (LLM) generates quiz questions from the transcript and POSTs the result to \texttt{POST /webhooks/ai-pipeline/quiz}. The system creates a \texttt{Quiz} aggregate in \texttt{DRAFT} status linked to the lesson.
    \item The instructor reviews the generated quiz and publishes it (see UC-06).
  \end{enumerate} \\
  \textbf{Postconditions} & The video is playable via HLS. A transcript is available. A draft quiz is created for instructor review. \\
  \textbf{Exceptions} &
  \begin{itemize}
    \item \textbf{E01: Transcoding failure.} If Mux reports \texttt{video.asset.errored}, the \texttt{VideoAsset} transitions to \texttt{ERRORED} and the instructor is notified.
    \item \textbf{E02: Transcription timeout.} If the AI pipeline does not respond within 30 minutes, the \texttt{TranscriptionJob} transitions to \texttt{TIMEOUT} and the instructor is notified to re-submit.
  \end{itemize} \\
  \bottomrule
\end{longtable}

% UC-05
\subsection{UC-05: Student Takes Quiz Attempt}

\begin{longtable}{p{3cm}p{10cm}}
  \toprule
  \textbf{Element} & \textbf{Description} \\
  \midrule
  \endhead
  \midrule
  \endfoot
  \textbf{ID} & UC-05 \\
  \textbf{Name} & Student Takes Quiz Attempt \\
  \textbf{Actors} & Authenticated student \\
  \textbf{Preconditions} & The student is enrolled in the course. The quiz is in \texttt{PUBLISHED} status and linked to a lesson. \\
  \textbf{Main Flow} &
  \begin{enumerate}
    \item The student requests the quiz via \texttt{GET /quizzes/:id}. The system returns the quiz questions (without answer keys) and the student's previous attempts, if any.
    \item The student submits answers via \texttt{POST /quizzes/:id/attempts}. The system creates a \texttt{QuizAttempt} aggregate in \texttt{SUBMITTED} status.
    \item For multiple-choice and true/false questions, the system grades automatically using the stored answer keys. For short-answer questions, the system publishes a \texttt{ShortAnswerGradeRequestedEvent} through the outbox.
    \item The AI pipeline grades short-answer questions asynchronously and delivers results via webhook. The \texttt{QuizAttempt} aggregate tracks per-question grading status (\texttt{PENDING}, \texttt{AUTO\_GRADED}, \texttt{AI\_GRADED}).
    \item When all questions are graded, the \texttt{QuizAttempt} transitions to \texttt{GRADED} and the total score and per-question feedback are computed.
    \item The student retrieves the graded attempt via \texttt{GET /attempts/:id} and views the score, correct/incorrect markings, and instructor feedback for short-answer questions.
  \end{enumerate} \\
  \textbf{Postconditions} & The quiz attempt is graded. The score is recorded in the student's progress. \\
  \textbf{Exceptions} &
  \begin{itemize}
    \item \textbf{E01: Attempt limit exceeded.} If the quiz has an attempt limit (e.g., 3 maximum) and the student has exhausted it, the system returns HTTP 403 Forbidden.
    \item \textbf{E02: Short-answer grading timeout.} If AI grading does not complete within 30 minutes, the attempt escalates to \texttt{REVIEW\_REQUIRED} for manual instructor intervention.
  \end{itemize} \\
  \bottomrule
\end{longtable}

% UC-06
\subsection{UC-06: Instructor Reviews AI-Generated Quiz}

\begin{longtable}{p{3cm}p{10cm}}
  \toprule
  \textbf{Element} & \textbf{Description} \\
  \midrule
  \endhead
  \midrule
  \endfoot
  \textbf{ID} & UC-06 \\
  \textbf{Name} & Instructor Reviews AI-Generated Quiz \\
  \textbf{Actors} & Authenticated instructor \\
  \textbf{Preconditions} & The AI pipeline has generated a draft quiz attached to a lesson. The quiz is in \texttt{DRAFT} status. \\
  \textbf{Main Flow} &
  \begin{enumerate}
    \item The instructor navigates to the lesson and sees a notification that an AI-generated quiz draft is available for review.
    \item The instructor requests the quiz via \texttt{GET /quizzes/:id} (with instructor permissions, which includes the answer keys).
    \item The instructor reviews each question, edits the text, changes the question type, adds or removes options, and adjusts the correct answer.
    \item The instructor publishes the quiz via \texttt{POST /quizzes/:id/publish}.
    \item The system validates that the quiz has at least one question and transitions it to \texttt{PUBLISHED} status.
  \end{enumerate} \\
  \textbf{Postconditions} & The quiz is available to students enrolled in the course. \\
  \textbf{Exceptions} &
  \begin{itemize}
    \item \textbf{E01: No questions.} If the instructor attempts to publish a quiz with zero questions, the system returns HTTP 422 Unprocessable Entity.
    \item \textbf{E02: Unauthorised.} If the authenticated user is not the course instructor or an admin, the system returns HTTP 403 Forbidden.
  \end{itemize} \\
  \bottomrule
\end{longtable}

% UC-07
\subsection{UC-07: Student Watches Video with Progress Tracking}

\begin{longtable}{p{3cm}p{10cm}}
  \toprule
  \textbf{Element} & \textbf{Description} \\
  \midrule
  \endhead
  \midrule
  \endfoot
  \textbf{ID} & UC-07 \\
  \textbf{Name} & Student Watches Video with Progress Tracking \\
  \textbf{Actors} & Authenticated student, Mux Video CDN \\
  \textbf{Preconditions} & The student is enrolled in the course. The lesson has a video asset in \texttt{READY} status. \\
  \textbf{Main Flow} &
  \begin{enumerate}
    \item The student requests the lesson via \texttt{GET /lessons/:id}. The system verifies enrollment and returns the lesson content and a signed Mux playback URL.
    \item The student's frontend loads the HLS stream using HLS.js and begins playback.
    \item The frontend sends periodic progress updates to \texttt{PATCH /lessons/:id/progress} with the current playback position (in seconds) and the percentage watched.
    \item When the student reaches 90\% of the video duration (configurable threshold), the frontend sends a completion event. The system creates or updates a \texttt{LessonProgress} record in \texttt{COMPLETED} status.
    \item The system publishes a \texttt{LessonCompleted} event, which triggers streak computation (via \texttt{StreakService}) and updates the \texttt{CourseProgressView}.
  \end{enumerate} \\
  \textbf{Postconditions} & The lesson is marked as completed in the student's progress. The streak counter is updated. If all lessons in the course are completed, the course is marked as completed. \\
  \textbf{Exceptions} &
  \begin{itemize}
    \item \textbf{E01: Enrollment not active.} If the enrollment is \texttt{CANCELLED} or \texttt{PENDING\_PAYMENT}, the system denies access and returns HTTP 403 Forbidden.
    \item \textbf{E02: Video not ready.} If the video asset is still processing, the system returns the lesson content without the video URL and a \texttt{Retry-After} header.
  \end{itemize} \\
  \bottomrule
\end{longtable}

% UC-08
\subsection{UC-08: Discussion Thread Creation and Instructor Reply}

\begin{longtable}{p{3cm}p{10cm}}
  \toprule
  \textbf{Element} & \textbf{Description} \\
  \midrule
  \endhead
  \midrule
  \endfoot
  \textbf{ID} & UC-08 \\
  \textbf{Name} & Discussion Thread Creation and Instructor Reply \\
  \textbf{Actors} & Authenticated student (thread creator), authenticated instructor (reply) \\
  \textbf{Preconditions} & The student is enrolled in the course. The course has discussion enabled. \\
  \textbf{Main Flow} &
  \begin{enumerate}
    \item The student creates a discussion thread via \texttt{POST /courses/:id/discussions} with a title and body. The system creates a \texttt{Thread} aggregate in \texttt{OPEN} status and publishes a \texttt{ThreadCreated} event.
    \item Other enrolled students can view the thread and post replies via \texttt{POST /discussions/:id/replies}.
    \item The instructor views the thread and posts a reply that is marked with the \texttt{IS\_INSTRUCTOR} flag. The system fires \texttt{InstructorReplied} and sends an email notification to the thread creator.
    \item The student can mark the thread as resolved via \texttt{PATCH /discussions/:id/resolve}, which transitions the thread to \texttt{RESOLVED} status.
  \end{enumerate} \\
  \textbf{Postconditions} & The thread is created and visible to all enrolled students. Participants can post replies. \\
  \textbf{Exceptions} &
  \begin{itemize}
    \item \textbf{E01: Rate limit exceeded.} If the student exceeds the maximum number of threads per hour (configurable), the system returns HTTP 429 Too Many Requests.
    \item \textbf{E02: Course not enrolled.} If the student is not enrolled in the course, the system returns HTTP 403 Forbidden.
  \end{itemize} \\
  \bottomrule
\end{longtable}

% UC-09
\subsection{UC-09: Certificate Generation After Course Completion}

\begin{longtable}{p{3cm}p{10cm}}
  \toprule
  \textbf{Element} & \textbf{Description} \\
  \midrule
  \endhead
  \midrule
  \endfoot
  \textbf{ID} & UC-09 \\
  \textbf{Name} & Certificate Generation After Course Completion \\
  \textbf{Actors} & Authenticated student \\
  \textbf{Preconditions} & The student has completed all lessons in the course (progress = 100\%). The course has certificate generation enabled. \\
  \textbf{Main Flow} &
  \begin{enumerate}
    \item When the last lesson in the course is marked as completed, the \texttt{CourseProgressView} detects that all lessons are completed and publishes a \texttt{CourseCompleted} event.
    \item The Certificates context handles the event and creates a \texttt{Certificate} aggregate in \texttt{PENDING} status.
    \item A background job generates a PDF certificate using a templating library (PDFKit or Puppeteer). The PDF includes the student's name, course title, completion date, and a verification QR code.
    \item The PDF is uploaded to AWS S3 with server-side encryption. The \texttt{Certificate} aggregate transitions to \texttt{ISSUED} status with an S3 key reference.
    \item The student can download the certificate via \texttt{GET /certificates/:id/download}, which generates a presigned S3 URL with 1-hour expiration.
    \item The student can share the certificate verification link: \texttt{https://nexalearn.app/verify/<certificate\_id>}, which displays the certificate details without requiring authentication.
  \end{enumerate} \\
  \textbf{Postconditions} & A verifiable PDF certificate is generated and stored. The student can download and share the certificate. \\
  \textbf{Exceptions} &
  \begin{itemize}
    \item \textbf{E01: Certificate already issued.} If the student requests a duplicate certificate, the system returns the existing certificate instead of generating a new one.
    \item \textbf{E02: S3 upload failure.} If the PDF upload to S3 fails, the \texttt{Certificate} aggregate remains in \texttt{PENDING} and a retry is scheduled via the outbox.
  \end{itemize} \\
  \bottomrule
\end{longtable}

% UC-10
\subsection{UC-10: Instructor Views Revenue Analytics Dashboard}

\begin{longtable}{p{3cm}p{10cm}}
  \toprule
  \textbf{Element} & \textbf{Description} \\
  \midrule
  \endhead
  \midrule
  \endfoot
  \textbf{ID} & UC-10 \\
  \textbf{Name} & Instructor Views Revenue Analytics Dashboard \\
  \textbf{Actors} & Authenticated instructor \\
  \textbf{Preconditions} & The instructor is logged in and has at least one published course with completed enrollments. \\
  \textbf{Main Flow} &
  \begin{enumerate}
    \item The instructor requests the analytics dashboard via \texttt{GET /analytics/instructor/dashboard}. The request includes query parameters for the time range (\texttt{?from=...\&to=...}).
    \item The Instructor context queries read models in the Progress and Payments contexts to compute:
      \begin{itemize}
        \item Total revenue for the period (sum of succeeded payment intents).
        \item Number of new enrollments per course per day.
        \item Average quiz score across all attempts.
        \item Course completion rate (students who completed / students enrolled).
        \item Top-performing courses by revenue and enrollment count.
      \end{itemize}
    \item The system returns the computed metrics as a JSON response. The frontend renders charts and tables from this data.
    \item The instructor can filter by course and date range and export the data as CSV.
  \end{enumerate} \\
  \textbf{Postconditions} & The instructor can view and export revenue and engagement analytics. \\
  \textbf{Exceptions} &
  \begin{itemize}
    \item \textbf{E01: No data.} If no enrollments or payments exist for the selected period, the system returns empty metrics with a descriptive message.
    \item \textbf{E02: Access denied.} If the instructor attempts to access analytics for a course they do not own, the system filters the data to include only owned courses.
  \end{itemize} \\
  \bottomrule
\end{longtable}
